% Options for packages loaded elsewhere
\PassOptionsToPackage{unicode}{hyperref}
\PassOptionsToPackage{hyphens}{url}
\documentclass[
  12pt,
]{ctexart}
\usepackage{xcolor}
\usepackage{amsmath,amssymb}
\setcounter{secnumdepth}{-\maxdimen} % remove section numbering
\usepackage{iftex}
\ifPDFTeX
  \usepackage[T1]{fontenc}
  \usepackage[utf8]{inputenc}
  \usepackage{textcomp} % provide euro and other symbols
\else % if luatex or xetex
  \usepackage{unicode-math} % this also loads fontspec
  \defaultfontfeatures{Scale=MatchLowercase}
  \defaultfontfeatures[\rmfamily]{Ligatures=TeX,Scale=1}
\fi
\usepackage{lmodern}
\ifPDFTeX\else
  % xetex/luatex font selection
\fi
% Use upquote if available, for straight quotes in verbatim environments
\IfFileExists{upquote.sty}{\usepackage{upquote}}{}
\IfFileExists{microtype.sty}{% use microtype if available
  \usepackage[]{microtype}
  \UseMicrotypeSet[protrusion]{basicmath} % disable protrusion for tt fonts
}{}
\makeatletter
\@ifundefined{KOMAClassName}{% if non-KOMA class
  \IfFileExists{parskip.sty}{%
    \usepackage{parskip}
  }{% else
    \setlength{\parindent}{0pt}
    \setlength{\parskip}{6pt plus 2pt minus 1pt}}
}{% if KOMA class
  \KOMAoptions{parskip=half}}
\makeatother
\usepackage{longtable,booktabs,array}
\usepackage{calc} % for calculating minipage widths
% Correct order of tables after \paragraph or \subparagraph
\usepackage{etoolbox}
\makeatletter
\patchcmd\longtable{\par}{\if@noskipsec\mbox{}\fi\par}{}{}
\makeatother
% Allow footnotes in longtable head/foot
\IfFileExists{footnotehyper.sty}{\usepackage{footnotehyper}}{\usepackage{footnote}}
\makesavenoteenv{longtable}
\setlength{\emergencystretch}{3em} % prevent overfull lines
\providecommand{\tightlist}{%
  \setlength{\itemsep}{0pt}\setlength{\parskip}{0pt}}
\usepackage{bookmark}
\IfFileExists{xurl.sty}{\usepackage{xurl}}{} % add URL line breaks if available
\urlstyle{same}
\hypersetup{
  hidelinks,
  pdfcreator={LaTeX via pandoc}}

\author{SCX}
\date{2026}

\begin{document}

\section{SCX
8篇定理论文数学完整性审查报告}\label{scx-8ux7bc7ux5b9aux7406ux8bbaux6587ux6570ux5b66ux5b8cux6574ux6027ux5ba1ux67e5ux62a5ux544a}

\subsection{0. SCX公理体系基准
(Baseline)}\label{scxux516cux7406ux4f53ux7cfbux57faux51c6-baseline}

审查前先厘清核心公理：

\begin{itemize}
\tightlist
\item
  \textbf{Theorem 1} (S1): 多专家一致性噪声检测 --- 在A1-A6假设下，F1 ≥
  1 - (1/η)Σρ\_s·exp(-2MΔ\_s²)，指数收敛。
\item
  \textbf{Theorem 2} (S2): 弱特征信息论边界 --- F1\_SCX ≤ F1\_base +
  C\_F·√(δ/2)，δ=I(φ(X);S)。
\item
  \textbf{Theorem 3} (S3): 噪声-难度不可区分性 ---
  从观测数据无法区分标签噪声与内在难度。这是SCX的''不确定性原理''。
\item
  \textbf{Theorem 4} (S4): 精确常数极小极大最优性 ---
  Bahadur-Rao大偏差精确常数。
\item
  \textbf{Cercis Score}: S(x) = Q(x) + η·N(x)，质量+新颖性分解。
\item
  \textbf{Situs算子}: 将分布编码为度量测度空间 (X, d\_P, μ\_P)。
\item
  \textbf{Spring SE-1}: 自演化门控，Robbins-Monro收敛保证。
\end{itemize}

\begin{center}\rule{0.5\linewidth}{0.5pt}\end{center}

\subsection{1. 逐篇审查}\label{ux9010ux7bc7ux5ba1ux67e5}

\subsubsection{1.1 Theorem 5 --- Active Learning (Cercis Optimal
Sampling)}\label{theorem-5-active-learning-cercis-optimal-sampling}

\textbf{定理陈述完整性}: ★★★★☆ (4/5) - Theorem 5.1 (Optimal
Distribution): 完整。Gibbs形式p\emph{(x) ∝
π(x)exp(αS(x)/λ)，有唯一性证明。 - Theorem 5.2 (Optimal Rate):
陈述完整，封闭解ρ} =
η/(1+η)·Σw\_k·V{[}Γ\_k{]}/Σw\_k·Γ̄\_k。但关键推导步骤有\textbf{严重逻辑跳跃}（见下文）。
- Theorem 5.3 (Consistency): 完整。√n-一致性，Delta方法。

\textbf{证明框架的逻辑跳跃}: ★★★☆☆ (3/5) --- \textbf{存在严重问题} -
Lemma 1 (Gibbs解) 的推导是标准的变分推理，没有问题。 -
\textbf{致命跳跃在Theorem 5.2的证明第414-421行}: 从ρ* =
η·Σw\_kΓ̄\_k/(η·Σw\_kΓ̄\_k + Σw\_kV{[}Γ\_k{]})跳到最终形式ρ* =
η/(1+η)·Σw\_kV{[}Γ\_k{]}/Σw\_kΓ̄\_k时，声称使用了恒等式 Σw\_kΓ̄\_k =
((1+η)/η)·Σw\_kV{[}Γ\_k{]}。这个恒等式被标注为''在变分目标的最优点处成立，来自等边际收益与η-缩放新颖性惩罚的边际成本相等''。\textbf{这是循环论证}------用ρ\emph{的定义来推导ρ}的封闭形式。实际上从代数看：
- 从Φ'(ρ)=0得 ρ = Σw\_kΓ̄\_k / (Σw\_kΓ̄\_k + η\^{}\{-1\}Σw\_kV{[}Γ\_k{]})
- 这个形式已经是封闭解，不需要额外的恒等式 -
论文强行转换到η/(1+η)形式时引入了一个假设的结构关系，破坏了封闭性 -
\textbf{Assumption 1 (Cercis-Gain Affinity)}: Γ̄(x) = α·S(x) +
β。这是整个推导的基石，但仅以''直觉''和''Taylor展开''论证，缺乏严格证明。这是一个未经证实的经验假设。

\textbf{与SCX公理体系的一致性}: ★★★★☆ (4/5) - 正确使用了Cercis Score
S=Q+ηN的定义 - η→0回收纯不确定性采样，η→∞回收纯多样性采样，与SCX框架一致
- 但''η/(1+η)``的饱和因子并未出现在原始SCX框架中------这是新引入的结构

\textbf{从现有公理可严格推导}: - Theorem 5.1 (Gibbs分布): ✓
从变分原理严格推导（无需SCX公理） - Theorem 5.3 (一致性): ✓
从大数定律+Delta方法严格推导

\textbf{需要新假设}: - Theorem 5.2 (采样率): ✗ 需要Assumption 1
(Cercis-Gain Affinity) + 额外的变分结构关系 -
风险调节函数Φ(ρ)的三项分解（期望收益、边际递减、新颖性成本×方差）是\textbf{ad
hoc构造}，不是从Cercis公理推导的

\textbf{与其他定理的关系}:
独立。被AR-Theorem的detectability结果引用（``T5最优主动学习策略''），但AR仅将T5作为黑箱工具使用。

\begin{center}\rule{0.5\linewidth}{0.5pt}\end{center}

\subsubsection{1.2 Theorem 6 --- Protocol Game (Audit Protocol
Adoption)}\label{theorem-6-protocol-game-audit-protocol-adoption}

\textbf{定理陈述完整性}: ★★★★☆ (4/5) - Theorem 6.1 (Regret bound):
完整。 - Theorem 6.2 (Nash existence): 完整，标准Glicksberg不动点论证。
- Theorem 6.3 (Uniqueness): 完整，显式γ\_crit封闭形式。 - Proposition
6.4 (Distortion): 完整。

\textbf{证明框架的逻辑跳跃}: ★★★★☆ (4/5) -
存在性证明是标准的------Glicksberg定理对连续收益的混合策略博弈直接适用。
- 唯一性证明的收缩映射论证正确，但存在一处\textbf{跳跃}:
影响力函数φ\_P对早期采纳者数量的敏感性被bound为√(2logN₁/(N₀+1)) +
σ\_ξ/√N₀，这个bound假设UCB置信半径的导数可以用有限差分近似。这个近似在大N₀下渐近成立，但有限样本下需要额外论证。
- Lemma
``影响力单调性''的证明仅依赖直觉论证（``增加n\_P减小标准误差''），缺乏对UCB索引随机性的形式化处理。

\textbf{与SCX公理体系的一致性}: ★★★☆☆ (3/5) -
论文声称''Quality(Q)对应协议质量μ\_P''，``Novelty(N)对应UCB探索奖励''。这个映射是\textbf{类比性的而非推导性的}------协议采纳博弈并非Cercis框架的自然延伸，而是独立构建的博弈论模型，然后事后与SCX类比。
-
γ\_crit包含UCB探索项√(2logN₁/(N₀+1))，这与Cercis的η参数没有直接数学联系，仅在概念层面平行。

\textbf{从现有公理可严格推导}: - Theorem 6.1 (Regret): ✓ 标准UCB分析 -
Theorem 6.2 (Existence): ✓ 标准博弈论

\textbf{需要新假设}: - Theorem 6.3 (Uniqueness): 需要Gumbel噪声假设 +
影响力敏感性bound（UCB有限差分近似） -
整个模型设定（顺序博弈、Gumbel偏好冲击、高斯先验）与SCX公理无关

\textbf{与其他定理的关系}:
与CD-Theorem有间接联系（信息外部性），与TS-Theorem的锚点概念平行。但\textbf{基本上是一篇独立论文}，挂靠SCX品牌而非从中推导。

\begin{center}\rule{0.5\linewidth}{0.5pt}\end{center}

\subsubsection{1.3 Theorem 7 --- Cross-Domain Partition
Preservation}\label{theorem-7-cross-domain-partition-preservation}

\textbf{定理陈述完整性}: ★★★★★ (5/5) - Theorem 7.1 (Main bound):
完整，三部分组成清晰。 - Theorem 7.2 (Tightness):
完整，构造了紧确性序列。 - Corollary 7.3 (最优粒度): 完整。

\textbf{证明框架的逻辑跳跃}: ★★★★☆ (4/5) - 五步证明结构合理。 -
\textbf{Step 3有潜在问题}: 声明''条件密度是L\_k-Lipschitz的 ⇒
熵泛函关于Wasserstein度量是L\_k-Lipschitz的''。熵不是一般Lipschitz函数，需要额外条件（如对数Sobolev不等式或有界支持上的条件密度有界）。Lemma
S2.2 (Entropy Lipschitz Property)
给出H(Y\textbar X=x)是(L/σ)-Lipschitz的论证，但这假设了条件方差有界且条件分布接近高斯------这是一个\textbf{重要但未明确声明的假设}。
- ``条件化不会增加平均传输成本''的论证（Step 4）是正确的标准结果。

\textbf{与SCX公理体系的一致性}: ★★★★★ (5/5) -
直接从Situs条件互信息I(Y;P\textbar S)出发 -
使用Situs编码和Wasserstein距离，完全在SCX框架内 -
明确引用了T7跨域传输定理

\textbf{从现有公理可严格推导}: - 主体bound: 部分可推导（需要Assumption 1
+ 熵Lipschitz条件） - 渐近紧确性(Tightness): ✓ 构造性证明

\textbf{需要新假设}: - Assumption 1 (Bounded Conditional Densities with
Lipschitz property) - Lemma S2.2的熵Lipschitz论证需要高斯性或等价条件

\textbf{与其他定理的关系}:
被CD-Theorem的因果中继定理引用为''T7跨域保真度''。是整个8篇中\textbf{数学上最自洽}的论文之一。

\begin{center}\rule{0.5\linewidth}{0.5pt}\end{center}

\subsubsection{1.4 CD-Theorem --- Causal Discovery (Noise vs
Confounding)}\label{cd-theorem-causal-discovery-noise-vs-confounding}

\textbf{定理陈述完整性}: ★★★★☆ (4/5) - CD-Theorem 1 (Strong
Inseparability): 完整。 - CD-Theorem 2 (Weak Separability):
陈述完整但\textbf{条件性极强}。 - CD-Theorem 3 (Causal Relay): 完整。

\textbf{证明框架的逻辑跳跃}: ★★★☆☆ (3/5) - \textbf{Strong
Inseparability的证明跳跃}: Step 1声称S(x) = S₀(x) + η·δ\_ε(x) +
β·δ\_z(x)，然后断言当η⊥z时，联合分布因子化使δ\_z不可区分于有效噪声。这是对Theorem
3的直接类比推广，但\textbf{论证过于简略}------需要证明(S, ∇S,
H\_S)的联合分布确实在两种世界中相同，而不仅仅是S的边际分布。Theorem
3的原始证明是通过显式构造两个世界的联合分布等价，CD-Theorem没有做这种显式构造。
- \textbf{Curl统计量的问题}: τ =
(1/\textbar X\textbar)∫‖∇×∇̂S‖²dx。这个统计量在计算上严重依赖Situs流形上的旋度定义。论文没有给出Situs流形上的旋度算子的具体构造，只是模糊引用了''SCX框架中的Situs流形''。从几何角度看，旋度需要联络结构（connection），而Situs编码仅提供度量------\textbf{从度量到旋度需要额外数学结构}（Levi-Civita联络在一般度量空间上不一定存在）。
- \textbf{Causal Relay}:
链式推导合理，但g(ε,δ)函数中的交互项C·ε·δ没有给出C的定义或bound方法。

\textbf{与SCX公理体系的一致性}: ★★★★☆ (4/5) - 直接扩展Theorem
3至因果发现环境 -
Situs骨架保持假设是对Situs算子的重大扩展，但论文诚实标注为开放问题 -
Causal Relay正确引用了T7

\textbf{从现有公理可严格推导}: - Strong Inseparability: 声称是Theorem
3的直接推论，但证明过于简略

\textbf{需要新假设}: - Situs骨架保持假设（核心开放问题） -
Situs流形上旋度算子的存在性 -
所有三个定理都有诚实标注（{[}Rigorous{]}/{[}Conditionally
Rigorous{]}/{[}Open{]}）

\textbf{与其他定理的关系}:
与AR-Theorem（梯度异常伪装为因果信号）和HC-Theorem（人类认知打破对称性）形成三角交叉引用。

\begin{center}\rule{0.5\linewidth}{0.5pt}\end{center}

\subsubsection{1.5 FA-Theorem --- Federated Audit (ZK
Multi-Party)}\label{fa-theorem-federated-audit-zk-multi-party}

\textbf{定理陈述完整性}: ★★★★☆ (4/5) - FA-Theorem 1 (ε-Equivalence):
完整，四个误差源。 - FA-Theorem 2 (Info Lower Bound): 完整。 - ZK
construction sketch: 完整但仅为概要。

\textbf{证明框架的逻辑跳跃}: ★★★☆☆ (3/5) - \textbf{核心未解决问题}:
Situs同态复合算子⊕的存在性。没有它，ε-Equivalence定理仅有条件意义。论文诚实承认这是开放问题。
- \textbf{I-MMSE gap的推导} (Step 2 of Lower Bound):
使用了Guo-Shamai-Verdú的I-MMSE关系，但该关系适用于高斯信道，而这里的Cercis
Score是非高斯的。论证需要推广或有额外辩护。 - \textbf{``最坏构造''}
(Step 3):
声称''对手可调整S̃\_k以最大化偏差''，但没有给出调整策略的具体形式。需要证明存在一个篡改使得E{[}T\_fed{]}与E{[}T\_full{]}的偏差达到声称的下界。

\textbf{与SCX公理体系的一致性}: ★★★★☆ (4/5) - 正确使用Cercis
Score和Situs编码 - Audit Sword原则是SCX框架的延伸 -
下界直接推导自Theorem 3

\textbf{从现有公理可严格推导}: - Info Lower Bound (Theorem 2): ✓
从Theorem 3 + 数据处理不等式推导

\textbf{需要新假设}: - Situs同态复合算子⊕（核心开放问题） -
ZK电路构造的可行性（密码学工程问题，非数学问题）

\textbf{与其他定理的关系}:
被CD-Theorem和AR-Theorem引用为''审计剑''原则的联邦化扩展。\textbf{是最具雄心的扩展论文}，但也是最依赖未解决数学问题的论文（Situs同态性）。

\begin{center}\rule{0.5\linewidth}{0.5pt}\end{center}

\subsubsection{1.6 AR-Theorem --- Adversarial Robustness (Cercis
Taxonomy)}\label{ar-theorem-adversarial-robustness-cercis-taxonomy}

\textbf{定理陈述完整性}: ★★★★☆ (4/5) - AR-Theorem 1 (Reduction):
完整，条件清晰。 - AR-Theorem 2 (Detectability): 完整，样本复杂度下界。
- Conjecture (Phase Transition): 陈述完整但仅为猜想。

\textbf{证明框架的逻辑跳跃}: ★★★☆☆ (3/5) -
\textbf{Reduction定理的关键跳跃}:
证明声称从\textbar\textbar∇S\textbar\textbar/S ≤
η（``来自Cercis算子定义''）得出A(x) ≤
η·ε。这个不等式\textbf{并非SCX框架中的标准结果}------Cercis算子定义S=Q+ηN，并没有直接约束\textbar\textbar∇S\textbar\textbar/S。这需要额外假设∇Q和∇N受限于Q和N的比例。
- \textbf{Taylor展开中的Hessian bound}:
\textbar\textbar H\_S\textbar\textbar\_op ≤ L(x) 是正确的（Lipschitz梯度
⇒ Hessian bounded），但仅在使用谱范数且函数两次可微时成立。 -
\textbf{Detectability定理的下界}: 使用Le Cam二点法推导χ²散度 =
Θ(γ²)，这一推导\textbf{极其简略}------需要显式计算梯度场分布在H₀与H₁之间的χ²散度，这依赖于对A(x)分布的具体假设。论文没有给出这个计算。
- \textbf{上界（T5可达性）}: 声称Theorem
5提供匹配上界，但T5的最优性是关于最大化I(S;query)的，而非关于检测对抗结构的速率最优性。这是\textbf{不同意义的最优性}------T5优化的是获取标签的样本效率，AR需要的是检测异常梯度结构的样本效率。

\textbf{与SCX公理体系的一致性}: ★★★☆☆ (3/5) -
三元分解S=Q+ηN+γA是对原始Cercis二元分解的重大扩展 -
是否在SCX框架内引入第三原语(A)是一个\textbf{公理层面的问题} -
梯度场诊断D\_adv的构造是新颖的，但在SCX框架中缺乏明确定义

\textbf{从现有公理可严格推导}: - Reduction的Lipschitz充分条件:
可推导，但使用了一个不在SCX公理中的不等式(\textbar\textbar∇S\textbar\textbar/S
≤ η)

\textbf{需要新假设}: - 三元分解S=Q+ηN+γA的合理性 -
\textbar\textbar∇S\textbar\textbar/S ≤ η（Cercis梯度-值比有界） -
对抗暴露A与新颖性N的不完全相关性（I(A;N\textbar S)\textgreater0） -
Situs流形上旋度算子（与CD-Theorem共享此问题）

\textbf{与其他定理的关系}:
明确引用CD（对抗信号与因果信号混淆）、HC（人类检测对抗样本）、TS（S\_crit漂移）。但在8篇中\textbf{数学基础最薄弱}------两个定理高度依赖未经验证的假设，核心猜想完全是开放的。

\begin{center}\rule{0.5\linewidth}{0.5pt}\end{center}

\subsubsection{1.7 TS-Theorem --- Temporal SCX (Drift \&
Self-Audit)}\label{ts-theorem-temporal-scx-drift-self-audit}

\textbf{定理陈述完整性}: ★★★★☆ (4/5) - TS-Theorem 1 (Drift
Detectability): 完整，样本量公式。 - TS-Theorem 2 (Drift Attribution):
完整，分无锚/有锚两种情况。 - TS-Theorem 3 (Sprint Self-Audit Limit):
完整，Lyapunov条件。

\textbf{证明框架的逻辑跳跃}: ★★★★☆ (4/5) - \textbf{Drift
Detectability证明是合理的}，但有一个隐含假设：Situs距离与η漂移之间的Lipschitz关系\textbar η\_\{t+1\}
- η\_t\textbar{} ≤ κ·d\_Situs +
\textbar ξ\_t\textbar。这个关系从T7的Lipschitz性质推导是合理的。 -
\textbf{方差公式(等式14)}包含σ²\_est项但没有定义或推导它。在H₀下，η的估计方差应能从Cercis
Score的分布推导，但论文直接断言了一个形式。 -
\textbf{归因不可识别性的构造}（Config A vs Config
B）有说服力，但仅针对两个配置，没有证明\textbf{所有}可能的分解都不可识别。需要更一般的不可识别性证明。
-
\textbf{Lyapunov收敛定理}是标准的Foster-Lyapunov论证，没有问题。发散条件的推导也是自然的。

\textbf{与SCX公理体系的一致性}: ★★★★☆ (4/5) -
明确将T7跨域传输定理时间化（Situs距离变为时间片间的Situs距离） -
Sprint自审计问题是SCX框架的自然时间扩展 -
锚点概念与HC-Theorem的锚点知识平行

\textbf{从现有公理可严格推导}: - Drift Detectability: ✓ 从T7 +
标准假设检验推导 - 归因不可识别性: ✓ 通过反例构造 - 锚定估计: ✓
从锚点定义直接推导

\textbf{需要新假设}: - Situs距离至η的Lipschitz常数κ -
锚点集合的存在性（实用但非理论需求） -
Lyapunov条件中的收缩率ρ（需要从Sprint更新规则推导）

\textbf{与其他定理的关系}: 与T7、HC共享锚点概念。收敛/发散分析与Spring
SE-1的Robbins-Monro收敛平行但不同------Spring是模型演化，TS是分布/审计参数演化。

\begin{center}\rule{0.5\linewidth}{0.5pt}\end{center}

\subsubsection{1.8 HC-Theorem --- Human-AI Collaborative
Audit}\label{hc-theorem-human-ai-collaborative-audit}

\textbf{定理陈述完整性}: ★★★★★ (5/5) - HC-Theorem 1 (Collaborative
Gain): 完整。 - HC-Theorem 2 (Budget Decay): 完整，√(B\_H/n)标度律。 -
HC-Theorem 3 (Absolute Boundary): 完整。

\textbf{证明框架的逻辑跳跃}: ★★★★☆ (4/5) - \textbf{Collaborative
Gain的逻辑}:
``如果人类判断携带超SCX信息，则协作审计严格优于纯AI''。这是一个条件蕴含，逻辑上正确，但\textbf{条件本身(I(J\_H;ε\textbar S)\textgreater0)是一个经验假设}。论文诚实承认这是开放问题。
- \textbf{Budget Decay的推导有微妙问题}: Step 1声称Ω\_collab = Ω\_AI +
(B\_H/n)·Δ\_I·η\_eff，但Step 2通过Cramér-Rao论证η\_eff ≤
c/√(B\_H/n)·1/Δ\_I，得到√(B\_H/n)标度。这个复合推导需要证明：基于Cercis
Score的最优选择确实产生√(B\_H/n)的效率损失。论文给出的论证过于简略------``最优选择引入样本间依赖，将有效样本量从B\_H降低到√(B\_H·n)''需要更严格的处理。
- \textbf{Absolute Boundary定理}: 这是Theorem
3的直接逻辑扩展，论证坚实。``完美噪声''的概念清晰：I(ε;所有可观测量\textbar X)=0
⇒ 绝对不可区分。

\textbf{与SCX公理体系的一致性}: ★★★★★ (5/5) - 直接扩展Theorem
3至任意认知系统 - ``多透镜''视角是对SCX框架的哲学深化 - 建立了从Theorem
3（AI-lens）到任意lens的连续统

\textbf{从现有公理可严格推导}: - Absolute Boundary (Theorem 3): ✓
Theorem 3的直接推论 - Budget Decay: 可从Theorem 3 +
信息论推导，但√(B\_H/n)标度的严格性需要更多工作

\textbf{需要新假设}: -
信息不对称条件I(J\_H;ε\textbar S)\textgreater0（经验假设） -
校准条件（人类不确定性随Cercis Score单调） -
效率因子η\_eff的上界推导需要更强的正则性条件

\textbf{与其他定理的关系}:
与CD（人类认知与SCX信号的交集）、AR（人类检测对抗样本）、TS（锚点）都有联系。是8篇中\textbf{哲学最深刻、对Theorem
3理解最透彻}的论文。

\begin{center}\rule{0.5\linewidth}{0.5pt}\end{center}

\subsection{2. 跨论文分析}\label{ux8de8ux8bbaux6587ux5206ux6790}

\subsubsection{2.1
定理重复或矛盾}\label{ux5b9aux7406ux91cdux590dux6216ux77dbux76fe}

\begin{itemize}
\tightlist
\item
  \textbf{重复}: 无直接重复。CD-Theorem的Strong Inseparability与Theorem
  3在精神上重叠但针对不同问题（噪声vs难度 vs 噪声vs混杂）。
\item
  \textbf{潜在矛盾}:

  \begin{itemize}
  \tightlist
  \item
    AR-Theorem的Detectability声称在∇S层次打破Theorem 3（``Theorem
    3的障碍在S层面适用，不在∇S层面''），而Theorem 3原始声明涵盖''Cercis
    Score及其导数''。如果Theorem
    3的原始证明确实涵盖∇S，那么AR的声称与Theorem
    3矛盾。\textbf{这需要在Theorem
    3的证明中检查是否真的涵盖了梯度信息。} Theorem 3原文说''observables
    accessible to an SCX auditor: the Cercis Score, its gradient, its
    Hessian''，结论是''No algorithm operating on these observables can
    tell the worlds apart''。AR-Theorem声称∇S中有Theorem
    3未穷尽的高阶结构（旋度）。\textbf{这是一个潜在的内部矛盾------要么Theorem
    3的声称被AR削弱，要么AR的声称无效。}
  \item
    T5的η/(1+η)因子与原始SCX的η参数语义不完全一致：原始框架中η是连续的探索-利用权衡，而T5中η通过η/(1+η)饱和因子影响采样率。这两个η是否应被视为同一参数？
  \end{itemize}
\end{itemize}

\subsubsection{\texorpdfstring{2.2 最弱论文: \textbf{AR-Theorem
(Adversarial
Robustness)}}{2.2 最弱论文: AR-Theorem (Adversarial Robustness)}}\label{ux6700ux5f31ux8bbaux6587-ar-theorem-adversarial-robustness}

\begin{itemize}
\tightlist
\item
  \textbf{理由}:

  \begin{enumerate}
  \def\labelenumi{\arabic{enumi}.}
  \tightlist
  \item
    核心结果之一（Reduction）依赖不等式\textbar\textbar∇S\textbar\textbar/S
    ≤ η，该不等式不在SCX公理中
  \item
    Detectability的χ²散度计算没有显式给出
  \item
    核心问题（Phase Transition Conjecture）完全是开放猜想
  \item
    三元分解S=Q+ηN+γA是对SCX公理的重大修改，但论文没有论证这种修改的合法性
  \item
    潜在的与Theorem 3矛盾
  \item
    Le Cam下界推导过于简略，可能包含隐藏假设
  \end{enumerate}
\item
  \textbf{诚实暴击}:
  这篇论文有一个好想法（对抗样本可能构成第三类），但数学执行不足。两个''定理''中的条件性过强，核心结论大多是''如果\ldots 那么\ldots``的条件语句，而条件本身未被满足或未经验证。实质上是一篇\textbf{定位论文+研究纲领}，不是已解决的定理。
\end{itemize}

\subsubsection{\texorpdfstring{2.3 最强论文: \textbf{Theorem 7
(Cross-Domain Partition
Preservation)}}{2.3 最强论文: Theorem 7 (Cross-Domain Partition Preservation)}}\label{ux6700ux5f3aux8bbaux6587-theorem-7-cross-domain-partition-preservation}

\begin{itemize}
\tightlist
\item
  \textbf{理由}:

  \begin{enumerate}
  \def\labelenumi{\arabic{enumi}.}
  \tightlist
  \item
    定理陈述最完整，每个符号都有明确定义
  \item
    证明五步分解清晰，引理辅助完善
  \item
    构造了渐近紧确性，这是最高标准的理论证明
  \item
    与SCX框架的一致性最高------直接从Situs条件互信息出发
  \item
    组件分解（几何×统计×不可约）提供可操作的实践指导
  \item
    扩展（p-Wasserstein、多目标域、自适应细化）展示了理论的丰富性
  \item
    \textbf{唯一的显著弱点}（熵Lipschitz需要高斯假设）在附录中有讨论
  \end{enumerate}
\item
  \textbf{次强候选}: HC-Theorem --- 对Theorem
  3的哲学深化极其出色，三层结构（AI-only → human-augmented → absolute
  limit）逻辑完美。Budget
  Decay的√(B\_H/n)标度是漂亮的信息论结果。唯一不足是核心假设依赖经验验证。
\end{itemize}

\begin{center}\rule{0.5\linewidth}{0.5pt}\end{center}

\subsection{3. 全局评估矩阵}\label{ux5168ux5c40ux8bc4ux4f30ux77e9ux9635}

\begin{longtable}[]{@{}llllll@{}}
\toprule\noalign{}
论文 & 完整性 & 逻辑严密性 & SCX一致性 & 新假设依赖 & 总体评分 \\
\midrule\noalign{}
\endhead
\bottomrule\noalign{}
\endlastfoot
T5-Active & ★★★★☆ & ★★★☆☆ & ★★★★☆ & 高 & \textbf{B+} \\
T6-Protocol & ★★★★☆ & ★★★★☆ & ★★★☆☆ & 高 & \textbf{B} \\
T7-CrossDomain & ★★★★★ & ★★★★☆ & ★★★★★ & 中 & \textbf{A} \\
CD-Causal & ★★★★☆ & ★★★☆☆ & ★★★★☆ & 高 & \textbf{B-} \\
FA-Federated & ★★★★☆ & ★★★☆☆ & ★★★★☆ & 极高 & \textbf{B-} \\
AR-Adversarial & ★★★★☆ & ★★☆☆☆ & ★★★☆☆ & 极高 & \textbf{C+} \\
TS-Temporal & ★★★★☆ & ★★★★☆ & ★★★★☆ & 中 & \textbf{B+} \\
HC-Human & ★★★★★ & ★★★★☆ & ★★★★★ & 中 & \textbf{A-} \\
\end{longtable}

\begin{center}\rule{0.5\linewidth}{0.5pt}\end{center}

\subsection{4. 关键发现总结}\label{ux5173ux952eux53d1ux73b0ux603bux7ed3}

\begin{enumerate}
\def\labelenumi{\arabic{enumi}.}
\item
  \textbf{T5的Theorem
  5.2证明存在代数/逻辑跳跃}------封闭形式的最终步骤引入了循环论证。但这个问题是可修复的（直接使用ρ
  = ΣΓ̄/(ΣΓ̄+η⁻¹ΣV)作为封闭形式即可，无需强行转换为η/(1+η)形式）。
\item
  \textbf{AR与Theorem 3存在潜在矛盾}------AR声称梯度场中有Theorem
  3未覆盖的信息，但Theorem 3明确声明涵盖梯度。需要CLARIFY。
\item
  \textbf{Situs同态复合算子(⊕)是FA-Theorem的核心瓶颈，也是整个联邦化SCX程序的阻塞点}。没有⊕，FA-Theorem降级为条件声明。
\item
  \textbf{CD-Theorem的旋度统计量依赖Situs流形上的联络结构}，这在纯度量空间中不天然存在。这是CD和AR共享的几何学问题。
\item
  \textbf{8篇论文中有4篇的核心结果在某种程度上是''条件定理''}（如果X成立，则Y成立），而X往往是开放问题或经验假设。这不是坏事------这正确定位了SCX研究的前沿------但读者需要知道哪些是已关闭的，哪些是开放的。
\item
  \textbf{最强的数学在T7}（信息论界+紧确性构造），\textbf{最深的哲学在HC}（多透镜理论+完美噪声的绝对边界），\textbf{最实用的在T5}（封闭形式采样率），\textbf{最大胆的在FA和AR}（但数学基础最不稳固）。
\end{enumerate}

\end{document}
